\documentclass[a4paper,14pt]{extarticle}

\input{C:/preamble.tex}

\usepackage{tcolorbox}

\usepackage{amsmath}
\usepackage{amssymb}
\usepackage{mathtools}

\begin{document}
Пусть $\Omega\subset\mathbb R^2$ --- ограниченная связная
область с липшицевой границей $\partial\Omega$.
Разобьём  $\partial\Omega$ на твёрдые стенки $\Gamma_D$,
вход $\Gamma_{\mathrm{L}}$ и выход $\Gamma_{\mathrm{R}}$.


Пусть $\mathbf n = (n_x,n_y)$ -- внешняя единичная нормаль
к $\partial \Omega$ . Введём функции
$
	p_{\mathrm{L}}~\text{и}~p_{\mathrm{R}}
$
на $\Gamma_{\mathrm{L}} ~\text{и}~\Gamma_{\mathrm{R}}$.




Тогда стационарная задача Стокса для переменной вязкости
$\mu$ 
\begin{equation}
	\begin{dcases}
		-\dfrac{\partial}{\partial x}
		\left(
			2\mu\dfrac{\partial v_x}{\partial x}
		\right)
		-
		\dfrac{\partial}{\partial y}
		\left[
			\mu\left(
				\dfrac{\partial v_x}{\partial y}
				+
				\dfrac{\partial v_y}{\partial x}
			\right)
		\right]
		+
		\dfrac{\partial p}{\partial x}
		= f_x, \\[6pt]
		-\dfrac{\partial}{\partial x}
		\left[
			\mu\left(
				\dfrac{\partial v_y}{\partial x}
				+
				\dfrac{\partial v_x}{\partial y}
			\right)
		\right]
		-
		\dfrac{\partial}{\partial y}
		\left(
			2\mu\dfrac{\partial v_y}{\partial y}
		\right)
		+
		\dfrac{\partial p}{\partial y}
		= f_y, \\[6pt]
		\dfrac{\partial v_x}{\partial x}
		+
		\dfrac{\partial v_y}{\partial y}
		= 0;
	\end{dcases}
	\label{eq1}
\end{equation}
здесь $v_x$ и $v_y$ --- компоненты скорости,
$p$ --- давление,  $f_x$ и $f_y$ --- компоненты заданной объёмной силы.
С граничными условиями прилипания на стенках и заданным
напряжением на входе и выходе
\begin{equation}
	\begin{gathered}
		v_x|_{\Gamma_D}=0,
		\qquad
		v_y|_{\Gamma_D}=0,
		\\[6pt]
		\mu\left(
			2\frac{\partial v_x}{\partial x}n_x
			+
			\left(
				\frac{\partial v_x}{\partial y}
				+
				\frac{\partial v_y}{\partial x}
			\right)n_y
		\right)
		-pn_x
		=
		-p_{\mathrm L}n_x,
		\quad \text{на }\Gamma_{\mathrm L},
		\\[6pt]
		\mu\left(
			2\frac{\partial v_x}{\partial x}n_x
			+
			\left(
				\frac{\partial v_x}{\partial y}
				+
				\frac{\partial v_y}{\partial x}
			\right)n_y
		\right)
		-pn_x
		=
		-p_{\mathrm R}n_x,
		\quad \text{на }\Gamma_{\mathrm R},
		\\[6pt]
		\mu\left(
			\left(
				\frac{\partial v_y}{\partial x}
				+
				\frac{\partial v_x}{\partial y}
			\right)n_x
			+
			2\frac{\partial v_y}{\partial y}n_y
		\right)
		-pn_y
		=
		-p_{\mathrm L}n_y,
		\quad \text{на }\Gamma_{\mathrm L},
		\\[6pt]
		\mu\left(
			\left(
				\frac{\partial v_y}{\partial x}
				+
				\frac{\partial v_x}{\partial y}
			\right)n_x
			+
			2\frac{\partial v_y}{\partial y}n_y
		\right)
		-pn_y
		=
		-p_{\mathrm R}n_y,
		\quad \text{на }\Gamma_{\mathrm R}.
	\end{gathered}
\end{equation}


Введём обозначения
\begin{equation}
	\begin{gathered}
	\mathbf u=(v_x,v_y)^T,
	\quad
	\mathbf f=(f_x,f_y)^T,
	\quad
	\Gamma_N=\Gamma_{\mathrm L}\cup\Gamma_{\mathrm R},
	\quad
	p_b=p_{\mathrm L}\cup p_{\mathrm R}. \\[6pt]
	\boldsymbol\sigma
	=
	-pI
	+
	\mu\left(
		\nabla\mathbf u+(\nabla\mathbf u)^T
	\right).
	\end{gathered}
\end{equation}

Тогда задача записывается в векторном виде:
\begin{equation}
	\begin{dcases}
		-\operatorname{div}\boldsymbol\sigma
		=
		\mathbf f,
		\\
		\operatorname{div}\mathbf u=0,
	\end{dcases}
	\qquad
	\begin{gathered}
		\left.\mathbf u\right|_{\Gamma_D}=0,
		\\
		\left.
		\boldsymbol\sigma\mathbf n
		\right|_{\Gamma_N}
		=
		-p_b\mathbf n.
	\end{gathered}
	\label{eq:stokes_vector}
\end{equation}
	
\section{Слабая форма}

Домножим первое и второе уравнения~\eqref{eq1}
на тестовые функции $\varphi$ и $\psi$, которые равны нулю
на $\Gamma_D$ и могут быть ненулевыми на $\Gamma_N$.
Третье уравнение домножим на тестовую функцию $q$.
Затем проинтегрируем уравнения по всей области.

Первое уравнение примет вид
\begin{equation}
	\begin{gathered}
		-\int_\Omega
		\frac{\partial}{\partial x}
		\left(
			2\mu\frac{\partial v_x}{\partial x}
		\right)
		\varphi\,d\Omega
		-
		\int_\Omega
		\frac{\partial}{\partial y}
		\left[
			\mu\left(
				\frac{\partial v_x}{\partial y}
				+
				\frac{\partial v_y}{\partial x}
			\right)
		\right]
		\varphi\,d\Omega
		+\\[6pt]
		+
		\int_\Omega
		\frac{\partial p}{\partial x}
		\varphi\,d\Omega
		=
		\int_\Omega f_x\varphi\,d\Omega.
	\end{gathered}
\end{equation}

Проинтегрируем слагаемые левой части по частям:
\begin{enumerate}
	\item Для слагаемых с производными скорости получаем
	\begin{equation*}
		\begin{gathered}
			-\int_\Omega
			\frac{\partial}{\partial x}
			\left(
				2\mu\frac{\partial v_x}{\partial x}
			\right)
			\varphi\,d\Omega
			-
			\int_\Omega
			\frac{\partial}{\partial y}
			\left[
				\mu\left(
					\frac{\partial v_x}{\partial y}
					+
					\frac{\partial v_y}{\partial x}
				\right)
			\right]
			\varphi\,d\Omega
			=\\[6pt]
			=
			\int_\Omega
			2\mu
			\frac{\partial v_x}{\partial x}
			\frac{\partial\varphi}{\partial x}\,d\Omega
			+
			\int_\Omega
			\mu\left(
				\frac{\partial v_x}{\partial y}
				+
				\frac{\partial v_y}{\partial x}
			\right)
			\frac{\partial\varphi}{\partial y}\,d\Omega
			-\\[6pt]
			-
			\int_{\partial\Omega}
			\mu\left[
				2\frac{\partial v_x}{\partial x}n_x
				+
				\left(
					\frac{\partial v_x}{\partial y}
					+
					\frac{\partial v_y}{\partial x}
				\right)n_y
			\right]
			\varphi\,d\Gamma.
		\end{gathered}
	\end{equation*}

	\item Для слагаемого с давлением:
	\begin{equation*}
		\int_\Omega
		\frac{\partial p}{\partial x}
		\varphi\,d\Omega
		=
		-\int_\Omega
		p\frac{\partial\varphi}{\partial x}\,d\Omega
		+
		\int_{\partial\Omega}
		pn_x\varphi\,d\Gamma.
	\end{equation*}
\end{enumerate}

Поскольку $\varphi=0$ на $\Gamma_D$, в граничных интегралах
остаётся только участок $\Gamma_N$. Объединяя граничные
слагаемые и перенося их в правую часть, получаем
\begin{equation}
	\begin{gathered}
		\int_\Omega
		2\mu
		\frac{\partial v_x}{\partial x}
		\frac{\partial\varphi}{\partial x}\,d\Omega
		+
		\int_\Omega
		\mu\left(
			\frac{\partial v_x}{\partial y}
			+
			\frac{\partial v_y}{\partial x}
		\right)
		\frac{\partial\varphi}{\partial y}\,d\Omega
		-\\[6pt]
		-
		\int_\Omega
		p\frac{\partial\varphi}{\partial x}\,d\Omega
		=
		\int_\Omega f_x\varphi\,d\Omega
		+\\[6pt]
		+
		\int_{\Gamma_N}
		\left\{
			\mu\left[
				2\frac{\partial v_x}{\partial x}n_x
				+
				\left(
					\frac{\partial v_x}{\partial y}
					+
					\frac{\partial v_y}{\partial x}
				\right)n_y
			\right]
			-pn_x
		\right\}
		\varphi\,d\Gamma.
	\end{gathered}
\end{equation}

Подставляя граничное условие и аналогично преобразуя
второе уравнение, получаем слабую форму:
\begin{equation}
	\begin{dcases}
		\int_\Omega
		2\mu
		\frac{\partial v_x}{\partial x}
		\frac{\partial\varphi}{\partial x}\,d\Omega
		+
		\int_\Omega
		\mu\left(
			\frac{\partial v_x}{\partial y}
			+
			\frac{\partial v_y}{\partial x}
		\right)
		\frac{\partial\varphi}{\partial y}\,d\Omega
		\\
		\quad
		-
		\int_\Omega
		p\frac{\partial\varphi}{\partial x}\,d\Omega
		=
		\int_\Omega f_x\varphi\,d\Omega
		-
		\int_{\Gamma_N}
		p_b n_x\varphi\,d\Gamma;
		\\[6pt]
		\int_\Omega
		\mu\left(
			\frac{\partial v_y}{\partial x}
			+
			\frac{\partial v_x}{\partial y}
		\right)
		\frac{\partial\psi}{\partial x}\,d\Omega
		+
		\int_\Omega
		2\mu
		\frac{\partial v_y}{\partial y}
		\frac{\partial\psi}{\partial y}\,d\Omega
		\\
		\quad
		-
		\int_\Omega
		p\frac{\partial\psi}{\partial y}\,d\Omega
		=
		\int_\Omega f_y\psi\,d\Omega
		-
		\int_{\Gamma_N}
		p_b n_y\psi\,d\Gamma;
		\\[6pt]
		\int_\Omega
		q\operatorname{div}\mathbf u\,d\Omega=0.
	\end{dcases}
	\label{eq:weak}
\end{equation}

Здесь $v_x,v_y\in H^1(\Omega)$, $p\in L^2(\Omega)$,
причём $v_x=v_y=0$ на $\Gamma_D$. Равенства должны
выполняться для всех $\varphi,\psi\in H^1(\Omega)$,
равных нулю на $\Gamma_D$, и всех $q\in L^2(\Omega)$.
Граничное условие для скорости понимается в смысле следа.

\section{Разбиение области и тестовые функции}

Разобьём область $\Omega$ на $N_\tau$ подобластей
$\tau_\alpha$ с непересекающимися внутренностями:
\begin{equation}
	\overline{\Omega}
	=
	\bigcup_{\alpha=1}^{N_\tau}
	\overline{\tau_\alpha}.
\end{equation}

Обозначим через $\gamma_\alpha$ часть границы Неймана,
принадлежащую границе подобласти $\tau_\alpha$:
\begin{equation}
	\gamma_\alpha
	=
	\partial\tau_\alpha \cap \Gamma_N.
\end{equation}

Предположим, что участки $\gamma_\alpha$ покрывают $\Gamma_N$
и не перекрываются по длине. Если подобласть $\tau_\alpha$
не соприкасается с $\Gamma_N$ по участку ненулевой длины,
то соответствующий граничный вклад равен нулю.

Интеграл по области представим суммой интегралов по подобластям,
а граничный интеграл --- суммой интегралов по участкам
$\gamma_\alpha$. Применяя это к первому уравнению слабой
формы~\eqref{eq:weak}, получаем
\begin{equation}
	\begin{gathered}
		\sum_{\alpha=1}^{N_\tau}
		\int_{\tau_\alpha}
		2\mu
		\frac{\partial v_x}{\partial x}
		\frac{\partial\varphi}{\partial x}
		\,d\tau_\alpha
		+
		\sum_{\alpha=1}^{N_\tau}
		\int_{\tau_\alpha}
		\mu\left(
			\frac{\partial v_x}{\partial y}
			+
			\frac{\partial v_y}{\partial x}
		\right)
		\frac{\partial\varphi}{\partial y}
		\,d\tau_\alpha
		-\\[6pt]
		-
		\sum_{\alpha=1}^{N_\tau}
		\int_{\tau_\alpha}
		p\frac{\partial\varphi}{\partial x}
		\,d\tau_\alpha
		=
		\sum_{\alpha=1}^{N_\tau}
		\int_{\tau_\alpha}
		f_x\varphi\,d\tau_\alpha
		-
		\sum_{\alpha=1}^{N_\tau}
		\int_{\gamma_\alpha}
		p_b n_x\varphi\,d\Gamma.
	\end{gathered}
	\label{eq:parts_first}
\end{equation}

Равенство выполняется для каждой допустимой глобальной
тестовой функции $\varphi$. Граничное слагаемое не исчезает,
так как тестовая функция на $\Gamma_N$ может быть ненулевой.

Выберем конечные наборы допустимых тестовых функций:
\begin{equation}
	\varphi_1,\ldots,\varphi_{m_\varphi},\qquad
	\psi_1,\ldots,\psi_{m_\psi},\qquad
	q_1,\ldots,q_{m_q}.
\end{equation}

Функции $\varphi_i$ и $\psi_i$ принадлежат $H^1(\Omega)$
и равны нулю на $\Gamma_D$, а $q_i\in L^2(\Omega)$.
Подставляя их в соответствующие уравнения, получаем систему
с конечным количеством строк:
\begin{equation}
	\begin{dcases}
		\sum_{\alpha=1}^{N_\tau}\int_{\tau_\alpha}
		2\mu
		\frac{\partial v_x}{\partial x}
		\frac{\partial\varphi_i}{\partial x}\,d\tau_\alpha
		+
		\sum_{\alpha=1}^{N_\tau}\int_{\tau_\alpha}
		\mu\left(
			\frac{\partial v_x}{\partial y}
			+
			\frac{\partial v_y}{\partial x}
		\right)
		\frac{\partial\varphi_i}{\partial y}\,d\tau_\alpha
		-
		\sum_{\alpha=1}^{N_\tau}\int_{\tau_\alpha}
		p\frac{\partial\varphi_i}{\partial x}\,d\tau_\alpha
		= \\
		=
		\sum_{\alpha=1}^{N_\tau}\int_{\tau_\alpha}
		f_x\varphi_i\,d\tau_\alpha
		-
		\sum_{\alpha=1}^{N_\tau}\int_{\gamma_\alpha}
		p_b n_x\varphi_i\,d\Gamma,
		\quad i=1,\ldots,m_\varphi;
		\\[6pt]
		\cdots\quad
		\text{(второе уравнение аналогично)}
		\quad i=1,\ldots,m_\psi;
		\\[6pt]
		\sum_{\alpha=1}^{N_\tau}\int_{\tau_\alpha}
		q_i\operatorname{div}\mathbf u\,d\tau_\alpha=0,
		\qquad i=1,\ldots,m_q.
	\end{dcases}
	\label{eq:parts_div}
\end{equation}

Получены равенства для выбранных тестовых функций. Искомые
функции здесь ещё не заменены конечномерными приближениями,
поэтому это пока не система относительно конечного числа
числовых коэффициентов.

Обозначим ограничения тестовых функций на подобласти:
\begin{equation}
	\Phi_{i,\alpha}
	=
	\varphi_i|_{\tau_\alpha},\qquad
	\Psi_{i,\alpha}
	=
	\psi_i|_{\tau_\alpha},\qquad
	Q_{i,\alpha}
	=
	q_i|_{\tau_\alpha}.
\end{equation}
	

Для первого набора ограничения можно расположить в матрицу
\begin{equation}
	\Phi
	=
	\begin{pmatrix}
		\Phi_{1,1}
		&
		\cdots
		&
		\Phi_{1,N_\tau}
		\\
		\vdots
		&
		\ddots
		&
		\vdots
		\\
		\Phi_{m_\varphi,1}
		&
		\cdots
		&
		\Phi_{m_\varphi,N_\tau}
	\end{pmatrix}.
\end{equation}

Аналогично задаются матрицы ограничений $\Psi$ и $\mathcal Q$.
Каждая строка содержит ограничения одной глобальной функции;
ограничения $\Phi_{i,\alpha}$ и $\Phi_{i,\beta}$ на общей
границе согласованы в смысле следов. Аналогичное замечание
относится к $\Psi$. Для функций $q_i\in L^2(\Omega)$
совпадение следов без дополнительных требований не предполагается.

Интеграл по подобласти $\tau_\alpha$ зависит только от
значений функций внутри этой подобласти. Поэтому вместо
функций определенных на всей области можно использовать их ограничения на $\tau_\alpha$: значения функций и их производных внутри подобласти при этом не изменяются.

Тогда первое уравнение принимает вид
\begin{equation}
	\begin{gathered}
		\sum_{\alpha=1}^{N_\tau}
		\int_{\tau_\alpha}
		2\mu
		\frac{\partial\left(v_x|_{\tau_\alpha}\right)}{\partial x}
		\frac{\partial\Phi_{i,\alpha}}{\partial x}
		\,d\tau_\alpha
		+
		\sum_{\alpha=1}^{N_\tau}
		\int_{\tau_\alpha}
		\mu\left(
			\frac{\partial\left(v_x|_{\tau_\alpha}\right)}{\partial y}
			+
			\frac{\partial\left(v_y|_{\tau_\alpha}\right)}{\partial x}
		\right)
		\frac{\partial\Phi_{i,\alpha}}{\partial y}
		\,d\tau_\alpha
		-\\[6pt]
		-
		\sum_{\alpha=1}^{N_\tau}
		\int_{\tau_\alpha}
		\left(p|_{\tau_\alpha}\right)
		\frac{\partial\Phi_{i,\alpha}}{\partial x}
		\,d\tau_\alpha
		=
		\sum_{\alpha=1}^{N_\tau}
		\int_{\tau_\alpha}
		\left(f_x|_{\tau_\alpha}\right)
		\Phi_{i,\alpha}\,d\tau_\alpha
		-
		\sum_{\alpha=1}^{N_\tau}
		\int_{\gamma_\alpha}
		p_b n_x\Phi_{i,\alpha}\,d\Gamma,
		\\[6pt]
		i=1,\ldots,m_\varphi.
	\end{gathered}
	\label{eq:restriction_x}
\end{equation}



\section{Аппроксимация решения}

В рассматриваемых пространствах можно выбрать полные
ортонормированные системы функций. Согласно теореме о базисе
сепарабельного гильбертова пространства, каждую функцию
такого пространства можно представить рядом по соответствующей
системе, сходящимся в норме этого пространства
(см.~\cite{hilbert_basis}).

Следовательно, отбрасывая все слагаемые ряда после некоторого
номера, можно получить конечную линейную комбинацию,
приближающую исходную функцию. При увеличении числа слагаемых
погрешность такого приближения стремится к нулю в норме
соответствующего пространства.

Поэтому численные решения будем искать в виде конечных линейных комбинаций базисных функций.
Для скорости будем выбирать базисные функции из
$H^1(\Omega)$ с нулевым следом на $\Gamma_D$, а для давления
--- из $L^2(\Omega)$. Поэтому конечные линейные комбинации
этих функций также принадлежат соответствующим пространствам,
а приближённая скорость удовлетворяет однородному граничному
условию на $\Gamma_D$.

\begin{equation}
	v_x \approx
	v_x^*
	=
	\sum_{j=1}^{N_x}\omega_{x,j}\lambda_{x,j},
	\quad
	v_y\approx
	v_y^*
	=
	\sum_{j=1}^{N_y}\omega_{y,j}\lambda_{y,j},
	\quad
	p \approx
	p^*
	=
	\sum_{k=1}^{N_p}\omega_{p,k}\lambda_{p,k},
\end{equation}
здесь $\omega_{x,j}$, $\omega_{y,j}$ и $\omega_{p,k}$ ---
неизвестные числовые коэффициенты.
Ограничивая эти разложения на подобласть $\tau_\alpha$,
получаем
\begin{equation}
		v_x^*|_{\tau_\alpha}
		=
		\sum_{j=1}^{N_x}
		\omega_{x,j}
		\left.\lambda_{x,j}\right|_{\tau_\alpha},
		\quad
		v_y^*|_{\tau_\alpha}
		=
		\sum_{j=1}^{N_y}
		\omega_{y,j}
		\left.\lambda_{y,j}\right|_{\tau_\alpha},\quad
		p^*|_{\tau_\alpha}
		=
		\sum_{k=1}^{N_p}
		\omega_{p,k}
		\left.\lambda_{p,k}\right|_{\tau_\alpha}.
	\label{eq:local_approximation}
\end{equation}

Теперь подставим эти разложения в первое уравнение, получим
\begin{equation}
	\begin{gathered}
		\sum_{\alpha=1}^{N_\tau}
		\int_{\tau_\alpha}
		2\mu
		\frac{\partial}{\partial x}
		\left(
			\sum_{j=1}^{N_x}
			\omega_{x,j}\lambda_{x,j,\alpha}
		\right)
		\frac{\partial\Phi_{i,\alpha}}{\partial x}
		\,d\tau_\alpha
		+\\[6pt]
		+
		\sum_{\alpha=1}^{N_\tau}
		\int_{\tau_\alpha}
		\mu
		\left[
			\frac{\partial}{\partial y}
			\left(
				\sum_{j=1}^{N_x}
				\omega_{x,j}\lambda_{x,j,\alpha}
			\right)
			+
			\frac{\partial}{\partial x}
			\left(
				\sum_{j=1}^{N_y}
				\omega_{y,j}\lambda_{y,j,\alpha}
			\right)
		\right]
		\frac{\partial\Phi_{i,\alpha}}{\partial y}
		\,d\tau_\alpha
		-\\[6pt]
		-
		\sum_{\alpha=1}^{N_\tau}
		\int_{\tau_\alpha}
		\left(
			\sum_{k=1}^{N_p}
			\omega_{p,k}\lambda_{p,k,\alpha}
		\right)
		\frac{\partial\Phi_{i,\alpha}}{\partial x}
		\,d\tau_\alpha
		=\\[6pt]
		=
		\sum_{\alpha=1}^{N_\tau}
		\int_{\tau_\alpha}
		f_x\Phi_{i,\alpha}\,d\tau_\alpha
		-
		\sum_{\alpha=1}^{N_\tau}
		\int_{\gamma_\alpha}
		p_b n_x\Phi_{i,\alpha}\,d\Gamma,
		\quad
		i=1,\ldots,m_\varphi.
	\end{gathered}
\end{equation}
Коэффициент $\omega_{x,j}$ соответствует глобальной базисной
функции $\lambda_{x,j}$ и поэтому остаётся одним и тем же
при рассмотрении её ограничений на разные подобласти.
В частности, если функция $\lambda_{x,j}$ ненулевая на
нескольких соседних подобластях, то во всех этих подобластях
при её ограничениях стоит один и тот же коэффициент
$\omega_{x,j}$. Аналогично это относится к коэффициентам
$\omega_{y,j}$ и $\omega_{p,k}$.

Поскольку коэффициенты $\omega_{x,j}$, $\omega_{y,j}$
и $\omega_{p,k}$ не зависят от координат, их можно вынести
за знаки интегралов. Используя линейность производной и
интеграла, а также меняя порядок конечных сумм, получаем
\begin{equation}
	\begin{gathered}
		\sum_{j=1}^{N_x}
		\omega_{x,j}
		\left[
			\sum_{\alpha=1}^{N_\tau}
			\int_{\tau_\alpha}
			\left(
				2\mu
				\frac{\partial\lambda_{x,j,\alpha}}{\partial x}
				\frac{\partial\Phi_{i,\alpha}}{\partial x}
				+
				\mu
				\frac{\partial\lambda_{x,j,\alpha}}{\partial y}
				\frac{\partial\Phi_{i,\alpha}}{\partial y}
			\right)
			\,d\tau_\alpha
		\right]
		+\\[6pt]
		+
		\sum_{j=1}^{N_y}
		\omega_{y,j}
		\left[
			\sum_{\alpha=1}^{N_\tau}
			\int_{\tau_\alpha}
			\mu
			\frac{\partial\lambda_{y,j,\alpha}}{\partial x}
			\frac{\partial\Phi_{i,\alpha}}{\partial y}
			\,d\tau_\alpha
		\right]
		-
		\sum_{k=1}^{N_p}
		\omega_{p,k}
		\left[
			\sum_{\alpha=1}^{N_\tau}
			\int_{\tau_\alpha}
			\lambda_{p,k,\alpha}
			\frac{\partial\Phi_{i,\alpha}}{\partial x}
			\,d\tau_\alpha
		\right]
		=\\[6pt]
		=
		\sum_{\alpha=1}^{N_\tau}
		\int_{\tau_\alpha}
		f_x\Phi_{i,\alpha}\,d\tau_\alpha
		-
		\sum_{\alpha=1}^{N_\tau}
		\int_{\gamma_\alpha}
		p_b n_x\Phi_{i,\alpha}\,d\Gamma,
		\quad
		i=1,\ldots,m_\varphi.
	\end{gathered}
\end{equation}
Множители коэффициентов нам известны, поэтому интегралы можно вычислить; неизвестными остаются
коэффициенты $\omega_{x,j}$, $\omega_{y,j}$ и $\omega_{p,k}$.

Введём обозначения
\begin{equation}
	\begin{gathered}
		(A_{xx})_{ij}
		=
		\sum_{\alpha=1}^{N_\tau}
		\int_{\tau_\alpha}
		\left(
			2\mu
			\frac{\partial\lambda_{x,j,\alpha}}{\partial x}
			\frac{\partial\Phi_{i,\alpha}}{\partial x}
			+
			\mu
			\frac{\partial\lambda_{x,j,\alpha}}{\partial y}
			\frac{\partial\Phi_{i,\alpha}}{\partial y}
		\right)
		\,d\tau_\alpha,
		\\[6pt]
		(A_{xy})_{ij}
		=
		\sum_{\alpha=1}^{N_\tau}
		\int_{\tau_\alpha}
		\mu
		\frac{\partial\lambda_{y,j,\alpha}}{\partial x}
		\frac{\partial\Phi_{i,\alpha}}{\partial y}
		\,d\tau_\alpha,
		\\[6pt]
		(B_x)_{ik}
		=
		\sum_{\alpha=1}^{N_\tau}
		\int_{\tau_\alpha}
		\lambda_{p,k,\alpha}
		\frac{\partial\Phi_{i,\alpha}}{\partial x}
		\,d\tau_\alpha,
		\\[6pt]
		(F_x)_i
		=
		\sum_{\alpha=1}^{N_\tau}
		\int_{\tau_\alpha}
		f_x\Phi_{i,\alpha}\,d\tau_\alpha
		-
		\sum_{\alpha=1}^{N_\tau}
		\int_{\gamma_\alpha}
		p_b n_x\Phi_{i,\alpha}\,d\Gamma,
		\\[6pt]
		\mathcal W_x
		=
		(\omega_{x,1},\dots,\omega_{x,N_x})^T,
		\quad
		\mathcal W_y
		=
		(\omega_{y,1},\dots,\omega_{y,N_y})^T,
		\quad
		\mathcal W_p
		=
		(\omega_{p,1},\dots,\omega_{p,N_p})^T.
	\end{gathered}
\end{equation}

Тогда получается система с $N_x+N_y+N_p$ столбцами
и $m_\varphi$ строками:
\begin{equation}
	A_{xx}\mathcal W_x
	+
	A_{xy}\mathcal W_y
	-
	B_x\mathcal W_p
	=
	F_x.
\end{equation}


Аналогично, вводя $A_{yx}$, $A_{yy}$, $B_y$ и $F_y$,
для второго уравнения получаем
\begin{equation}
	A_{yx}\mathcal W_x
	+
	A_{yy}\mathcal W_y
	-
	B_y\mathcal W_p
	=
	F_y.
\end{equation}

Для третьего уравнения ($i=1,\ldots,m_q$)
\begin{equation}
	\begin{gathered}
		\sum_{j=1}^{N_x}
		\omega_{x,j}
		\left(
			\sum_{\alpha=1}^{N_\tau}
			\int_{\tau_\alpha}
			Q_{i,\alpha}
			\frac{\partial\lambda_{x,j,\alpha}}{\partial x}
			\,d\tau_\alpha
		\right)
		+
		\sum_{j=1}^{N_y}
		\omega_{y,j}
		\left(
			\sum_{\alpha=1}^{N_\tau}
			\int_{\tau_\alpha}
			Q_{i,\alpha}
			\frac{\partial\lambda_{y,j,\alpha}}{\partial y}
			\,d\tau_\alpha
		\right)
		=0
	\end{gathered}
\end{equation}
вводим
\begin{equation}
		(C_x)_{ij}
		=
		\sum_{\alpha=1}^{N_\tau}
		\int_{\tau_\alpha}
		Q_{i,\alpha}
		\frac{\partial\lambda_{x,j,\alpha}}{\partial x}
		\,d\tau_\alpha,\quad
		(C_y)_{ij}
		=
		\sum_{\alpha=1}^{N_\tau}
		\int_{\tau_\alpha}
		Q_{i,\alpha}
		\frac{\partial\lambda_{y,j,\alpha}}{\partial y}
		\,d\tau_\alpha.
\end{equation}
тогда
\begin{equation}
	C_x\mathcal W_x+C_y\mathcal W_y=0.
\end{equation}

Полную систему можно записать с помощью блочной матрицы:
\begin{equation}
	\begin{pmatrix}
		A_{xx} & A_{xy} & -B_x\\
		A_{yx} & A_{yy} & -B_y\\
		C_x & C_y & 0
	\end{pmatrix}
	\begin{pmatrix}
		\mathcal W_x\\
		\mathcal W_y\\
		\mathcal W_p
	\end{pmatrix}
	=
	\begin{pmatrix}
		F_x\\
		F_y\\
		0
	\end{pmatrix}.
\end{equation}

Условия Неймана учтены в правых частях $F_x$ и $F_y$.
Условия Дирихле учтены при выборе базисных функций скорости:
они равны нулю на $\Gamma_D$, поэтому приближённая скорость
также равна нулю на этой части границы.

При реализации можно сначала собрать расширенную систему
для всех узлов скорости, включая граничные. Затем для
каждого узла на $\Gamma_D$ значения коэффициентов
$\omega_{x,j}$ и $\omega_{y,j}$ фиксируются равными нулю.
Для этого соответствующие строки и столбцы матрицы
обнуляются, на диагонали устанавливаются единицы, а
соответствующие элементы правой части приравниваются к нулю.

После этого полученная система готова к решению.




\section{Выбор базисных функций}

Для построения приближённого решения воспользуемся методом
Галёркина. Выберем тестовые функции совпадающими с базисными
функциями, используемыми для приближения соответствующих
искомых функций:
$
	\varphi_i=\lambda_{x,i},
$ 
$
\psi_i=\lambda_{y,i},
$
$
	q_i=\lambda_{p,i}.
$
Тогда количество тестовых функций совпадает с количеством
соответствующих базисных функций:
$
	m_\varphi=N_x,
$
$
	m_\psi=N_y,
$
$
	m_q=N_p.
$

Тогда возникают соотношения
\begin{equation}
		(C_x)_{ki}
		=
		\sum_{\alpha=1}^{N_\tau}
		\int_{\tau_\alpha}
		\lambda_{p,k,\alpha}
		\frac{\partial\lambda_{x,i,\alpha}}{\partial x}
		\,d\tau_\alpha
		=
		(B_x)_{ik},
		\quad \text{аналогично} \quad
		(C_y)_{ki}
		=
		(B_y)_{ik}.
\end{equation}

Таким образом, глобальную систему можно записать в виде
\begin{equation}
	\begin{pmatrix}
		A_{xx} & A_{xy} & -B_x\\
		A_{yx} & A_{yy} & -B_y\\
		B_x^T & B_y^T & 0
	\end{pmatrix}
	\begin{pmatrix}
		\mathcal W_x\\
		\mathcal W_y\\
		\mathcal W_p
	\end{pmatrix}
	=
	\begin{pmatrix}
		F_x\\
		F_y\\
		0
	\end{pmatrix}.
\end{equation}
Полученная матрица является квадратной и имеет размер
$(N_x+N_y+N_p)\times(N_x+N_y+N_p)$.

Для построения устойчивой дискретной системы необходимо
подходящим образом выбрать пространства приближений скорости
и давления. Одним из таких выборов является элемент
Тейлора--Худа: область разюивается на треугольники, на каждом скорость приближается полиномами второй степени, а давление --- полиномами первой
степени. Устойчивость соответствующей пары пространств,
в том числе выполнение дискретного условия inf-sup при
определённых условиях на сетку и граничные условия,
доказана в литературе~\cite{taylor_hood}.

Имеем $N_T$ треугольников $T_h$. Поскольку
$\left.p^*\right|_{T_h}\in P_1(T_h)$ для каждого
$h=1,\ldots,N_T$, то это ограничение однозначно задаётся
тремя коэффициентами $(c_1,c_2,c_3)$ в базисе $\{1,x,y\}$.
Следовательно, для задания базиса в $P_1(T_h)$ необходимы
три линейно независимые функции.

В качестве базиса можно использовать не только $\{1,x,y\}$,
но и другие три линейно независимые функции из $P_1(T_h)$.
Выберем базис, связанный с вершинами треугольника.


Для того чтобы ограничения $p^*|_{T_a}$ и $p^*|_{T_b}$
совпадали на общем ребре $T_a\cap T_b$, достаточно
совпадения их значений в двух различных точках этого ребра,
поскольку оба ограничения линейны.

Удобно выбрать базисные функции так, чтобы коэффициенты
разложения на треугольнике $\omega_{p,i}^{(h)}$ были равны
значениям приближённого давления в его вершинах $a_j = (x_j, y_j)$:
\begin{equation}
	p^*\left(a_j^{(h)}\right)
	=
	\omega_{p,j}^{(h)},
	\qquad j=1,2,3;
\end{equation}
подставляя разложение давления, получаем
\begin{equation}
	\sum_{i=1}^{3}\omega_{p,i}^{(h)}\lambda_{p,i}^{(h)}(a_j^{(h)})
	=\omega_{p,j}^{(h)};
\end{equation}
	отсюда получаем требования
\begin{equation}
		\lambda_{p,i}^{(h)}(a_j^{(h)})=\delta_{ij},\quad \text{и} \quad
		\lambda_{p,1}^{(h)}+	\lambda_{p,2}^{(h)}+	\lambda_{p,3}^{(h)} = 1.
\end{equation}

\begin{tcolorbox}
		Рассмотрим эталонный треугольник
		$\widehat{T}\subset\mathbb{R}^2_{\xi,\eta}$
		с вершинами
		$M_1=(0,0)$, $M_2=(1,0)$, $M_3=(0,1)$.
		На нем базисные функции
		пространства $P_1(\widehat{T})$,
		удовлетворяющие условиям~\eqref{eq:qwe},
		имеют вид
		\vspace{-8pt}
		\begin{equation}
			\phi_1(\xi,\eta)=1-\xi-\eta,\qquad
			\phi_2(\xi,\eta)=\xi,\qquad
			\phi_3(\xi,\eta)=\eta.
		\end{equation}
	\end{tcolorbox}
	
	Для построения базисных функций на произвольном
	треугольнике $T_h$ сетки можно воспользоваться
	следующим подходом:
	\begin{enumerate}
		\item ввести  $\widehat T$;
		
		\item для треугольника с вершинами $(x_1,y_1)$,
		$(x_2,y_2)$, $(x_3,y_3)$ построить аффинное отображение
		$G_h:T_h\to\widehat T$, определяемое системой
		\begin{equation}
			\begin{pmatrix}
				x_2-x_1 & x_3-x_1\\
				y_2-y_1 & y_3-y_1
			\end{pmatrix}
			\begin{pmatrix}
				\xi\\
				\eta
			\end{pmatrix}
			=
			\begin{pmatrix}
				x-x_1\\
				y-y_1
			\end{pmatrix},
			\qquad G_h(x,y)=(\xi,\eta);
		\end{equation}
		
		
		\item на $\widehat T$ задать $\phi_1,\phi_2,\phi_3$;
		
		\item с помощью обратного отображения
		$G_h^{-1}$ перенести $\phi_1,\phi_2,\phi_3$
		на  $T_h$,
		получив 
		$\lambda_{p,i}^{(h)}(x,y)$:
	\begin{equation}
		\lambda^{(h)}_{p,i}(x,y) = \phi_i(G_h(x,y)),\quad (x,y) \in T_h.
	\end{equation}
	\end{enumerate}
	В результате они будут выглядеть следующим образом
	\begin{equation}
	\begin{gathered}
	D^{(h)} = (x_2 - x_1)(y_3 - y_1) - (x_3 - x_1)(y_2 - y_1),
	\\
	\lambda_{p,1}^{(h)}(x,y) =
	\frac{(x_2 - x)(y_3 - y) - (x_3 - x)(y_2 - y)}{D^{(h)}},
	\\
	\lambda_{p,2}^{(h)}(x,y) =
	\frac{(y_3 - y_1)(x - x_1) - (x_3 - x_1)(y - y_1)}{D^{(h)}},
	\\
	\lambda_{p,3}^{(h)}(x,y) =
	\frac{(x_2 - x_1)(y - y_1) - (y_2 - y_1)(x - x_1)}{D^{(h)}},
	\end{gathered}
	\quad (x,y) \in T_h.
	\end{equation}
	%Запишем каждую базисную функцию в виде
	%\[
	%\lambda_{p,i}^{(h)}(x,y)
	%=
	%c_{i,1}+c_{i,2}x+c_{i,3}y,
	%\qquad i=1,2,3.
	%\]
Тогда коэффициенты определяются следующим образом:
	\begin{equation}
		\begin{pmatrix}
			c_{1,1} & c_{1,2} & c_{1,3}\\
			c_{2,1} & c_{2,2} & c_{2,3}\\
			c_{3,1} & c_{3,2} & c_{3,3}
		\end{pmatrix}
		=
		\frac{1}{D^{(h)}}
		\begin{pmatrix}
			x_2y_3-x_3y_2 & y_2-y_3 & x_3-x_2\\
			x_3y_1-x_1y_3 & y_3-y_1 & x_1-x_3\\
			x_1y_2-x_2y_1 & y_1-y_2 & x_2-x_1
		\end{pmatrix}.
	\end{equation}

	Для приближения каждой компоненты скорости на треугольнике
	$T_h$ используется пространство $P_2(T_h)$ размерности 6.
	Поэтому выберем шесть узлов: три вершины и середины трёх
	рёбер. Коэффициенты разложения, как и для давления, свяжем
	со значениями приближённого решения в этих узлах.

	Ограничение квадратичной функции на ребро однозначно
	определяется значениями в трёх точках. Поэтому для
	согласования скоростей на общем ребре достаточно совпадения
	их значений в двух его вершинах и середине. Получаем похожие условия
	\begin{equation}
	\lambda_{u,i}^{(h)}\left(a_j^{(h)}\right)
	=
	\delta_{ij},
	\quad i,j=1,\ldots,6 \quad \text{и} \quad
	\lambda_1 + \cdots + \lambda_6 = 1.
	\end{equation}
	Подставляя координаты шести узлов в выражение для
$\lambda_{u,i}^{(h)}$, получаем систему относительно
коэффициентов $z_{i,1},\ldots,z_{i,6}$:
\begin{equation}
	z_{i,1}+z_{i,2}x_j+z_{i,3}y_j
	+z_{i,4}x_j^2+z_{i,5}x_jy_j+z_{i,6}y_j^2
	=
	\delta_{ij},
	\qquad j=1,\ldots,6.
\end{equation}
Для невырожденного треугольника эта система имеет
единственное решение. Чтобы найти его, воспользуемся
ранее построенными базисными функциями давления и выразим
через них базисные функции скорости:
\begin{equation}
	\begin{gathered}
		\lambda_{u,1}^{(h)}
		=
		\lambda_{p,1}^{(h)}
		\left(2\lambda_{p,1}^{(h)}-1\right),
		\quad
		\lambda_{u,2}^{(h)}
		=
		\lambda_{p,2}^{(h)}
		\left(2\lambda_{p,2}^{(h)}-1\right),
		\quad
		\lambda_{u,3}^{(h)}
		=
		\lambda_{p,3}^{(h)}
		\left(2\lambda_{p,3}^{(h)}-1\right),
		\\[6pt]
		\lambda_{u,4}^{(h)}
		=
		4\lambda_{p,1}^{(h)}\lambda_{p,2}^{(h)},
		\quad
		\lambda_{u,5}^{(h)}
		=
		4\lambda_{p,2}^{(h)}\lambda_{p,3}^{(h)},
		\quad
		\lambda_{u,6}^{(h)}
		=
		4\lambda_{p,3}^{(h)}\lambda_{p,1}^{(h)}.
	\end{gathered}
\end{equation}
Первые три функции соответствуют вершинам треугольника,
а остальные --- серединам рёбер $a_1$--$a_2$, $a_2$--$a_3$ и
$a_3$--$a_1$ соответственно. Эти функции удовлетворяют нужным
условия.

Подставляя в них ранее найденные линейные функции давления
и приводя подобные слагаемые, получаем коэффициенты $z_{ij}$.


Для функций, соответствующих вершинам, имеем
\begin{equation}
	\begin{gathered}
		z_{i,1}=2c_{i,1}^2-c_{i,1},
		\quad
		z_{i,2}=(4c_{i,1}-1)c_{i,2},
		\quad
		z_{i,3}=(4c_{i,1}-1)c_{i,3},
		\\[6pt]
		z_{i,4}=2c_{i,2}^2,
		\quad
		z_{i,5}=4c_{i,2}c_{i,3},
		\quad
		z_{i,6}=2c_{i,3}^2,
		\qquad i=1,2,3.
	\end{gathered}
\end{equation}
Для функций, соответствующих серединам рёбер, получаем
\begin{equation}
	\begin{gathered}
		z_{i,1}=4c_{r,1}c_{s,1},
		\quad
		z_{i,2}=4\left(c_{r,1}c_{s,2}+c_{s,1}c_{r,2}\right),
		\quad
		z_{i,3}=4\left(c_{r,1}c_{s,3}+c_{s,1}c_{r,3}\right),
		\\[6pt]
		z_{i,4}=4c_{r,2}c_{s,2},
		\quad
		z_{i,5}=4\left(c_{r,2}c_{s,3}+c_{s,2}c_{r,3}\right),
		\quad
		z_{i,6}=4c_{r,3}c_{s,3},
		\\[6pt]
		(i,r,s)\in
		\{(4,1,2),(5,2,3),(6,3,1)\}.
	\end{gathered}
\end{equation}



\newpage
\section{Сопряжённая задача Стокса--Дарси}

Текперь рассмотрим область свободного течения $\Omega_S$ и
пористую область $\Omega_D$. Их общую границу обозначим
через $\mathcal I_{\partial}$.

Скорость и давление в области Стокса переобозначим через
$\mathbf v_S$ и $p_S$, а скорость фильтрации и давление
в пористой области  $\mathbf v_D$ и $p_D$.


В области $\Omega_S$ все так же сохраним
$\Gamma_{SD}$ и $\Gamma_{SN} = \Gamma_{\mathrm L}\cup\Gamma_{\mathrm R}$ для условий Дирихле и Неймана. 

Внешнюю границу $\Omega_D$ для граничного условия непроницаемости обозначим $\Gamma_{DN}$.



\begin{tcolorbox}
	Уравнение Дарси и усовие непроницаемости выглядят:
	\begin{equation}
		\begin{dcases}
			v_{D,x}
			=
			-\dfrac{K}{\mu}
			\dfrac{\partial p_D}{\partial x},
			\\[6pt]
			v_{D,y}
			=
			-\dfrac{K}{\mu}
			\dfrac{\partial p_D}{\partial y},
			\\[6pt]
			\dfrac{\partial v_{D,x}}{\partial x}
			+
			\dfrac{\partial v_{D,y}}{\partial y}
			=
			0,
		\end{dcases}
		\quad
		\left.
		\left(
			v_{D,x}n_{D,x}
			+
			v_{D,y}n_{D,y}
		\right)
		\right|_{\Gamma_{DN}}
		=
		0;
	\end{equation}
	где $K$ --- проницаемость пористой среды, $\mu$ ---
	вязкость жидкости.

	В векторном виде, обозначив
	$\mathbf v_D=(v_{D,x},v_{D,y})^T$ -- 
	\begin{equation}
		\begin{dcases}
			\mathbf v_D
			=
			-\dfrac{K}{\mu}\nabla p_D,
			\\[6pt]
			\operatorname{div}\mathbf v_D=0,
		\end{dcases}
		\qquad
		\left.
		(\mathbf v_D\cdot\mathbf n_D)
		\right|_{\Gamma_{DN}}
		=0.
	\end{equation}
\end{tcolorbox}


Объединим уравнения с заданными условиями на внешних
границах. Первые два уравнения рассматриваются в $\Omega_S$,
следующие два --- в $\Omega_D$:
\begin{equation}
	\begin{dcases}
		-\operatorname{div}\boldsymbol\sigma_S
		=
		\mathbf f,
		\\[6pt]
		\operatorname{div}\mathbf v_S=0,
		\\[6pt]
		\mathbf v_D
		=
		-\dfrac{K}{\mu}\nabla p_D,
		\\[6pt]
		\operatorname{div}\mathbf v_D=0;
	\end{dcases}
	\qquad
	\begin{dcases}
		\left.\mathbf v_S\right|_{\Gamma_{SD}}=0,
		\\[6pt]
		\left.
		\boldsymbol\sigma_S\mathbf n_S
		\right|_{\Gamma_{SN}}
		=
		-p_b\mathbf n_S,
		\\[6pt]
		\left.
		(\mathbf v_D\cdot\mathbf n_D)
		\right|_{\Gamma_{DN}}
		=0.
	\end{dcases}
	\label{eq:stokes_darcy_outer}
\end{equation}





Для замыкания задачи зададим условия сопряжения на общей
границе $\mathcal I_{\partial}$. Пусть $\mathbf n_S$ и
$\mathbf n_D$ --- внешние единичные нормали к  $\Omega_S$ и $\Omega_D$, а $\boldsymbol\tau$ ---
единичная касательная к $\mathcal I_{\partial}$. На
$\mathcal I_{\partial}$ нормали направлены противоположно:
\[
	\mathbf n_S=-\mathbf n_D.
\]

Даьше перейдем непосредственно к самими условиям.

\subsection{Условие непрерывности нормального потока}
Потребуем выполнения закона сохранения массы на общей
границе. Для произвольного участка
$I\subset\mathcal I_{\partial}$ его интегральная форма
имеет вид
\begin{equation}
	\int_I
	\left(
		\mathbf v_S\cdot\mathbf n_S
		+
		\mathbf v_D\cdot\mathbf n_D
	\right)d\Gamma
	=
	0.
\end{equation}

Поскольку участок $I$ выбран произвольно, получаем
локальную форму условия:
\begin{equation}
	\mathbf v_S\cdot\mathbf n_S
	+
	\mathbf v_D\cdot\mathbf n_D
	=
	0
	\quad\text{на}\quad\mathcal I_{\partial}.
	\label{eq:normal_flux}
\end{equation}

Поскольку $\mathbf n_D=-\mathbf n_S$, нормальные компоненты
скоростей по разные стороны границы совпадают:
\begin{equation}
	\mathbf v_S\cdot\mathbf n_S
	=
	\mathbf v_D\cdot\mathbf n_S.
\end{equation}

Если нормаль $\mathbf n_S$ образует угол $\theta$ с осью $x$,
то есть $\mathbf n_S=(\cos\theta,\sin\theta)^T$, условие
можно записать покомпонентно:
\begin{equation}
	v_{S,x}\cos\theta+v_{S,y}\sin\theta
	=
	v_{D,x}\cos\theta+v_{D,y}\sin\theta.
\end{equation}

Таким образом, скачок нормальной компоненты скорости через
границу исключается. Касательные компоненты по разные
стороны границы могут различаться: условие сохранения
массы само по себе их равенства не требует. В трёхмерном
случае это относится к компонентам скорости в касательной
плоскости. Их связь определяется отдельным условием
касательного скольжения.


\subsection{Баланс нормальных напряжений}
В приближении Стокса пренебрегаем инерционным переносом
импульса $\operatorname{div}
(\rho\mathbf v_S\otimes\mathbf v_S)$. Будем также считать,
что нормальное воздействие со стороны $\Omega_D$
определяется давлением $p_D$, а дополнительные нормальные
поверхностные силы на границе отсутствуют. Тогда для
произвольного участка $I\subset\mathcal I_{\partial}$
нормальную составляющую баланса импульса запишем в виде
\begin{equation}
	\int_I
	\left[
		(\boldsymbol\sigma_S\mathbf n_S)\cdot\mathbf n_S
		+
		p_D
	\right]d\Gamma
	=
	0;
\end{equation}
в локальной форме --
\begin{equation}
	\boxed{
		(\boldsymbol\sigma_S\mathbf n_S)\cdot\mathbf n_S
		=
		-p_D
		\quad\text{на}\quad\mathcal I_{\partial}.
	}
	\label{eq:normal_stress}
\end{equation}

Таким образом, давления $p_S$ и $p_D$ в общем случае
не обязаны совпадать: их разность связана с нормальной
составляющей вязкого напряжения в области Стокса.
Равенство $p_S=p_D$ получается, когда этот вязкий вклад
равен нулю.

Полученное условие относится только к нормальной
составляющей напряжения. Для касательной составляющей
используется отдельное условие скольжения.




\subsection{Баланс касательных напряжений (Условие Биверса-Джозефа-Саффмана)}

Обозначим через $\mathbf t_D$ силу, с которой свободная жидкость
действует на пористую среду, то есть на жидкость в порах и твёрдый
каркас, отнесённую к единице площади границы $\mathcal I_\partial$.

На произвольном участке $I\subset\mathcal I_\partial$
рассмотрим касательную составляющую баланса сил, предполагая,
что дополнительные поверхностные силы и накопление импульса
на границе отсутствуют:
\begin{equation}
	\int_I
	\left(
		\boldsymbol\sigma_S\mathbf n_S
		+
		\mathbf t_D
	\right)
	\cdot\boldsymbol\tau\,d\Gamma
	=
	0,
\end{equation}
где $\boldsymbol\tau$ --- единичный касательный вектор
к границе $\mathcal I_\partial$.

Поскольку равенство выполняется для произвольного участка
$I$, в локальной форме получаем
\begin{equation}
	(\boldsymbol\sigma_S\mathbf n_S)
	\cdot\boldsymbol\tau
	+
	\mathbf t_D\cdot\boldsymbol\tau
	=
	0
	\quad\text{на }\mathcal I_\partial.
	\label{eq:tangential_balance}
\end{equation}

Равенство \eqref{eq:tangential_balance} связывает касательное
напряжение $(\boldsymbol\sigma_S\mathbf n_S)
\cdot\boldsymbol\tau$, действующее на свободную жидкость,
и касательную силу $\mathbf t_D\cdot\boldsymbol\tau$,
действующую на пористую среду.

Для нормальной составляющей в п.~5.2 было принято, что
воздействие пористой среды определяется давлением $p_D$.
Закон Дарси сам по себе не определяет касательную силу,
поэтому для замыкания равенства
\eqref{eq:tangential_balance} необходимо задать
дополнительную связь между касательной силой и скоростью
течения.

Касательное взаимодействие свободной жидкости с пористой
средой связано, в частности, с трением о её твёрдый каркас.
Эксперименты Биверса и Джозефа показали, что жидкость может
проскальзывать вдоль поверхности пористой среды. Поэтому
для описания течения вдоль такой границы вместо условия
прилипания используют условие скольжения.

В рассматриваемой модели запишем условие
Биверса--Джозефа в виде
\begin{equation}
	(\boldsymbol\sigma_S\mathbf n_S)
	\cdot\boldsymbol\tau
	=
	-\frac{\alpha\mu}{\sqrt{K}}\,
	(\mathbf v_S-\mathbf v_D)\cdot\boldsymbol\tau
	\quad\text{на }\mathcal I_\partial,
	\label{eq:BJ}
\end{equation}
где $\alpha>0$ --- безразмерный эмпирический параметр,
зависящий от структуры пористой среды.

Следуя Саффману, предположим, что касательная составляющая
скорости фильтрации $\mathbf v_D\cdot\boldsymbol\tau$ мала
по сравнению с касательной составляющей скорости свободного
течения $\mathbf v_S\cdot\boldsymbol\tau$, и пренебрежём ею.
Тогда из \eqref{eq:BJ} и \eqref{eq:tangential_balance}
получаем
\begin{equation}
	\mathbf t_D\cdot\boldsymbol\tau
	=
	\beta\,\mathbf v_S\cdot\boldsymbol\tau,
	\qquad
	\beta=\frac{\alpha\mu}{\sqrt{K}}>0.
	\label{eq:friction_force}
\end{equation}

Подставляя это выражение в
\eqref{eq:tangential_balance}, получаем условие
Биверса--Джозефа--Саффмана:
\begin{equation}
	(\boldsymbol\sigma_S\mathbf n_S)
	\cdot\boldsymbol\tau
	=
	-\beta\,\mathbf v_S\cdot\boldsymbol\tau
	\quad\text{на }\mathcal I_\partial.
	\label{eq:BJS}
\end{equation}

Знак в \eqref{eq:BJS} означает, что касательное воздействие
пористой среды на свободную жидкость направлено против
касательной составляющей её скорости. Коэффициент $\beta$
характеризует сопротивление скольжению. Его размерность
равна $\mathrm{Па\cdot с/м}$, поэтому произведение
$\beta\,\mathbf v_S\cdot\boldsymbol\tau$ имеет размерность
напряжения.

Формально при $\beta\to 0$ условие \eqref{eq:BJS}
переходит в условие нулевого касательного напряжения.
При $\beta\to\infty$, если касательное напряжение остаётся
ограниченным, касательная составляющая скорости стремится
к нулю.

В трёхмерном случае условие \eqref{eq:BJS} записывается
для двух линейно независимых касательных векторов
$\boldsymbol\tau_1$ и $\boldsymbol\tau_2$ к границе
$\mathcal I_\partial$:
\begin{equation}
	(\boldsymbol\sigma_S\mathbf n_S)
	\cdot\boldsymbol\tau_i
	=
	-\beta\,\mathbf v_S\cdot\boldsymbol\tau_i,
	\qquad i=1,2.
\end{equation}

При принятых предположениях нормальная составляющая
векторного баланса сил
\[
	\boldsymbol\sigma_S\mathbf n_S+\mathbf t_D=\mathbf 0
\]
совпадает с условием \eqref{eq:normal_stress}: из него
следует $\mathbf t_D\cdot\mathbf n_S=p_D$. Таким образом,
п.~5.2 и 5.3 описывают нормальную и касательную проекции
баланса сил на границе.



\subsection{Итоговые вид задачи сопряжения}
Объединяя ранее написанное:
\begin{equation}
	\begin{dcases}
		-\operatorname{div}\boldsymbol\sigma_S
		=
		\mathbf f,
		\\[6pt]
		\operatorname{div}\mathbf v_S
		=
		0,
		\\[6pt]
		\mathbf v_D
		=
		-\dfrac{K}{\mu}\nabla p_D,
		\\[6pt]
		\operatorname{div}\mathbf v_D
		=
		0
	\end{dcases}
	;
	\quad
	\begin{dcases}
		\left.\mathbf v_S\right|_{\Gamma_{SD}}=0,
		\\[6pt]
		\left.
		\boldsymbol\sigma_S\mathbf n_S
		\right|_{\Gamma_{SN}}
		=
		-p_b\mathbf n_S,
		\\[6pt]
		\left.
		(\mathbf v_D\cdot\mathbf n_D)
		\right|_{\Gamma_{DN}}
		=0
	\end{dcases}
	; \quad
	\begin{dcases}
		\mathbf v_S\cdot\mathbf n_S
		+
		\mathbf v_D\cdot\mathbf n_D
		=
		0,
		\\[6pt]
		(\boldsymbol\sigma_S\mathbf n_S)
		\cdot\mathbf n_S
		=
		-p_D,
		\\[6pt]
		(\boldsymbol\sigma_S\mathbf n_S)
		\cdot\boldsymbol\tau
		=
		-\beta\,\mathbf v_S\cdot\boldsymbol\tau
	\end{dcases}
	\quad 
	\text{на }
	\mathcal{I}.
\end{equation}

\section{Слабая форма задачи сопряжения}
$\varphi$ и $\psi$ могут быть ненулевыми на $\mathcal I_\partial$. Тогда разница с \eqref{1.2} будет только в  первых двух уравнениях и несущественной, а именно в
\begin{equation}
	\begin{gathered}
		\int_{\partial\Omega_S}
		(\boldsymbol\sigma_S\mathbf n_S)\cdot
		\begin{pmatrix}
			\varphi\\
			\psi
		\end{pmatrix}
		\,d\Gamma
		=
		\int_{\Gamma_{SN}}
		(\boldsymbol\sigma_S\mathbf n_S)\cdot
		\begin{pmatrix}
			\varphi\\
			\psi
		\end{pmatrix}
		\,d\Gamma
		+
		\int_{\mathcal I_\partial}
		(\boldsymbol\sigma_S\mathbf n_S)\cdot
		\begin{pmatrix}
			\varphi\\
			\psi
		\end{pmatrix}
		\,d\Gamma.
	\end{gathered}
	\label{eq:stokes_boundary_split}
\end{equation}
Разложим
тестовые функции на нормальную и касательную:
\begin{equation}
	\begin{pmatrix}
		\varphi\\
		\psi
	\end{pmatrix}
	=
	\left(
		\begin{pmatrix}
			\varphi\\
			\psi
		\end{pmatrix}
		\cdot\mathbf n_S
	\right)\mathbf n_S
	+
	\left(
		\begin{pmatrix}
			\varphi\\
			\psi
		\end{pmatrix}
		\cdot\boldsymbol\tau
	\right)\boldsymbol\tau;
	\label{eq:test_normal_tangent}
\end{equation}
тогда
\begin{equation}
	\begin{gathered}
		\int_{\mathcal I_\partial}
		(\boldsymbol\sigma_S\mathbf n_S)\cdot
		\begin{pmatrix}
			\varphi\\
			\psi
		\end{pmatrix}
		\,d\Gamma
		=
		\\[6pt]
		=
		\int_{\mathcal I_\partial}
		\bigl(
			(\boldsymbol\sigma_S\mathbf n_S)
			\cdot\mathbf n_S
		\bigr)
		\left(
			\begin{pmatrix}
				\varphi\\
				\psi
			\end{pmatrix}
			\cdot\mathbf n_S
		\right)
		\,d\Gamma
		+
		\int_{\mathcal I_\partial}
		\bigl(
			(\boldsymbol\sigma_S\mathbf n_S)
			\cdot\boldsymbol\tau
		\bigr)
		\left(
			\begin{pmatrix}
				\varphi\\
				\psi
			\end{pmatrix}
			\cdot\boldsymbol\tau
		\right)
		\,d\Gamma.
	\end{gathered}
	\label{eq:interface_integral_split}
\end{equation}
Подставляя в \eqref{eq:interface_integral_split} \eqref{}
\begin{equation}
	\begin{gathered}
		\int_{\mathcal I_\partial}
		(\boldsymbol\sigma_S\mathbf n_S)\cdot
		\begin{pmatrix}
			\varphi\\
			\psi
		\end{pmatrix}
		\,d\Gamma
		=
		-
		\int_{\mathcal I_\partial}
		p_D
		\left(
			\begin{pmatrix}
				\varphi\\
				\psi
			\end{pmatrix}
			\cdot\mathbf n_S
		\right)
		\,d\Gamma
		-
		\int_{\mathcal I_\partial}
		\beta
		(\mathbf v_S\cdot\boldsymbol\tau)
		\left(
			\begin{pmatrix}
				\varphi\\
				\psi
			\end{pmatrix}
			\cdot\boldsymbol\tau
		\right)
		\,d\Gamma.
	\end{gathered}
	\label{eq:interface_integral_result}
\end{equation}
Переносем интгралы содержащие неизветсные $\mathbf v_S$ и $p_D$ в левую часть и запишем первые два уравнения в векторном виде полностью
\begin{equation}
	\begin{gathered}
		\int_{\Omega_S}
		\boldsymbol\sigma_S :
		\nabla
		\begin{pmatrix}
			\varphi\\
			\psi
		\end{pmatrix}
		\,d\Omega
		+
		\int_{\mathcal I_\partial}
		p_D
		\left(
			\begin{pmatrix}
				\varphi\\
				\psi
			\end{pmatrix}
			\cdot\mathbf n_S
		\right)
		\,d\Gamma
		+
		\int_{\mathcal I_\partial}
		\beta(\mathbf v_S\cdot\boldsymbol\tau)
		\left(
			\begin{pmatrix}
				\varphi\\
				\psi
			\end{pmatrix}
			\cdot\boldsymbol\tau
		\right)
		\,d\Gamma
		=
		\\[6pt]
		=
		\int_{\Omega_S}
		\mathbf f\cdot
		\begin{pmatrix}
			\varphi\\
			\psi
		\end{pmatrix}
		\,d\Omega
		-
		\int_{\Gamma_{SN}}
		p_b\mathbf n_S\cdot
		\begin{pmatrix}
			\varphi\\
			\psi
		\end{pmatrix}
		\,d\Gamma.
	\end{gathered}
	\label{eq:stokes_coupled_weak}
\end{equation}
Уравнение несжимаемости Стокса, как и в~(1.2), имеет вид
\begin{equation}
	\int_{\Omega_S}
	q\operatorname{div}\mathbf v_S\,d\Omega
	=
	0.
	\label{eq:stokes_coupled_div}
\end{equation}


Подставляя закон Дарси в уравнение несжимаемости в области
$\Omega_D$, получаем
\begin{equation}
	-\operatorname{div}
	\left(
		\frac{K}{\mu}\nabla p_D
	\right)
	=
	0
	\quad\text{в }\Omega_D.
	\label{eq:darcy_pressure}
\end{equation}
Умножим \eqref{eq:darcy_pressure} на тестовую функцию $g$
и проинтегрируем по области $\Omega_D$:
\begin{equation}
	-\int_{\Omega_D}
	\operatorname{div}
	\left(
		\frac{K}{\mu}\nabla p_D
	\right)
	g\,d\Omega
	=
	0.
	\label{eq:darcy_integral}
\end{equation}
Преобразуем подынтегральное выражение, используя правило
дифференцирования произведения:
\begin{equation}
	\begin{gathered}
		\operatorname{div}
		\left(
			g\frac{K}{\mu}\nabla p_D
		\right)
		=
		g\operatorname{div}
		\left(
			\frac{K}{\mu}\nabla p_D
		\right)
		+
		\frac{K}{\mu}\nabla p_D\cdot\nabla g.
	\end{gathered}
\end{equation}
Применяя формулу Гаусса,
получаем
\begin{equation}
	\begin{gathered}
		-\int_{\Omega_D}
		\operatorname{div}
		\left(
			\frac{K}{\mu}\nabla p_D
		\right)
		g\,d\Omega
		=
		\int_{\Omega_D}
		\frac{K}{\mu}\nabla p_D\cdot\nabla g\,d\Omega
		-
		\int_{\partial\Omega_D}
		\frac{K}{\mu}\nabla p_D\cdot\mathbf n_D
		\,g\,d\Gamma.
	\end{gathered}
	\label{eq:darcy_integration_by_parts}
\end{equation}

Используя закон Дарси, запишем
\begin{equation}
	-\frac{K}{\mu}\nabla p_D\cdot\mathbf n_D
	=
	\mathbf v_D\cdot\mathbf n_D.
\end{equation}
Тогда \eqref{eq:darcy_parts} принимает вид
\begin{equation}
	\int_{\Omega_D}
	\frac{K}{\mu}\nabla p_D\cdot\nabla g\,d\Omega
	+
	\int_{\partial\Omega_D}
	(\mathbf v_D\cdot\mathbf n_D)g\,d\Gamma
	=
	0.
	\label{eq:darcy_boundary_flux}
\end{equation}

Разбивая граничный интеграл и используя условия
непроницаемости и непрерывности потока, получаем
\begin{equation}
	\begin{gathered}
		\int_{\partial\Omega_D}
		(\mathbf v_D\cdot\mathbf n_D)g\,d\Gamma
		=
		\int_{\Gamma_{DN}}
		(\mathbf v_D\cdot\mathbf n_D)g\,d\Gamma
		+
		\int_{\mathcal I_\partial}
		(\mathbf v_D\cdot\mathbf n_D)g\,d\Gamma
		=
		-
		\int_{\mathcal I_\partial}
		(\mathbf v_S\cdot\mathbf n_S)g\,d\Gamma.
	\end{gathered}
\end{equation}

Таким образом, слабая форма уравнения Дарси имеет вид
\begin{equation}
	\int_{\Omega_D}
	\frac{K}{\mu}\nabla p_D\cdot\nabla g\,d\Omega
	-
	\int_{\mathcal I_\partial}
	(\mathbf v_S\cdot\mathbf n_S)g\,d\Gamma
	=
	0.
	\label{eq:darcy_weak_coupled}
\end{equation}


Собирая все уравнения получаем
\begin{equation}
	\begin{dcases}
		\int_{\Omega_S}
		\boldsymbol\sigma_S :
		\nabla
		\begin{pmatrix}
			\varphi\\
			\psi
		\end{pmatrix}
		\,d\Omega
		+
		\int_{\mathcal I_\partial}
		p_D
		\left(
			\begin{pmatrix}
				\varphi\\
				\psi
			\end{pmatrix}
			\cdot\mathbf n_S
		\right)
		\,d\Gamma
		\\
		+
		\int_{\mathcal I_\partial}
		\beta
		(\mathbf v_S\cdot\boldsymbol\tau)
		\left(
			\begin{pmatrix}
				\varphi\\
				\psi
			\end{pmatrix}
			\cdot\boldsymbol\tau
		\right)
		\,d\Gamma
		=
		\\
		=
		\int_{\Omega_S}
		\mathbf f\cdot
		\begin{pmatrix}
			\varphi\\
			\psi
		\end{pmatrix}
		\,d\Omega
		-
		\int_{\Gamma_{SN}}
		p_b\mathbf n_S\cdot
		\begin{pmatrix}
			\varphi\\
			\psi
		\end{pmatrix}
		\,d\Gamma,
		\\[6pt]
		\int_{\Omega_S}
		q\operatorname{div}\mathbf v_S\,d\Omega
		=
		0,
		\\[6pt]
		\int_{\Omega_D}
		\frac{K}{\mu}\nabla p_D\cdot\nabla g\,d\Omega
		-
		\int_{\mathcal I_\partial}
		(\mathbf v_S\cdot\mathbf n_S)g\,d\Gamma
		=
		0.
	\end{dcases}
	\label{eq:stokes_darcy_weak}
\end{equation}
для
$\mathbf v_S\in[H^1(\Omega_S)]^2$,
$p_S\in L^2(\Omega_S)$, 
$p_D\in H^1(\Omega_D)$.


\subsection*{По аналогии с разделом 2} 
Разобьём области $\Omega_S$ и
$\Omega_D$ на подобласти:
\begin{equation}
	\overline{\Omega}_S
	=
	\bigcup_{\alpha=1}^{N_S}
	\overline{\tau^S_\alpha},
	\qquad
	\overline{\Omega}_D
	=
	\bigcup_{\alpha=1}^{N_D}
	\overline{\tau^D_\alpha},
\end{equation}
которые пересекаются не более чем на линии.
Предположим, что разбиения согласованы на общей границе
$\mathcal I_\partial$: тогд ее также можно представить в виде
объединения участков $I_l$, которые пересекаются не более чем на точкку, каждый из них является
общим для одной подобласти Стокса и одной подобласти Дарси:
\begin{equation}
	\mathcal I_\partial
	=
	\bigcup_{l=1}^{N_I} I_l,
\end{equation} 

Добавим к наборам тестовых функций из раздела 2 также набор $g_1,\ldots,g_{m_g}\in H^1(\Omega_D)$,
обозначим $ G_{i,\alpha} = g_i|_{\tau^D_\alpha},~
i=1,\ldots,m_g,~ \alpha=1,\ldots,N_D. $

Для участков границы сопряжения введём обозначения
\begin{equation}
	\Phi^{\mathcal I_\partial}_{i,l}
	=
	\varphi_i|_{I_l},
	\qquad
	\Psi^{\mathcal I_\partial}_{i,l}
	=
	\psi_i|_{I_l},
	\qquad
	G^{\mathcal I_\partial}_{i,l}
	=
	g_i|_{I_l};
\end{equation}
$\varphi_i$ и $\psi_i$ рассматриваются со стороны
$\Omega_S$, а $g_i$ --- со стороны $\Omega_D$.

Для давления в области Стокса изменим обозначения, а для
давления в области Дарси введём разложение:
\begin{equation}
	p_S^*
	=
	\sum_{k=1}^{N_{p_S}}
	\omega_{S,k}\lambda_{S,k},
	\qquad
	p_D^*
	=
	\sum_{k=1}^{N_{p_D}}
	\omega_{D,k}\lambda_{D,k},
\end{equation}
где $\omega_{S,k}$ и $\omega_{D,k}$ --- неизвестные
числовые коэффициенты, а $\lambda_{S,k}$ и $\lambda_{D,k}$
--- базисные функции, определённые в $\Omega_S$ и
$\Omega_D$ соответственно.

Ограничивая разложения на подобласти, получаем
\begin{equation}
	\begin{gathered}
		p_S^*|_{\tau^S_\alpha}
		=
		\sum_{k=1}^{N_{p_S}}
		\omega_{S,k}
		\left(\lambda_{S,k}|_{\tau^S_\alpha}\right),
		\qquad
		p_D^*|_{\tau^D_\alpha}
		=
		\sum_{k=1}^{N_{p_D}}
		\omega_{D,k}
		\left(\lambda_{D,k}|_{\tau^D_\alpha}\right).
	\end{gathered}
\end{equation}

Для участков границы сопряжения $I_l$ имеем
\begin{equation}
	\begin{gathered}
		v_{S,x}^*|_{I_l}
		=
		\sum_{j=1}^{N_x}
		\omega_{x,j}
		\left(\lambda_{x,j}|_{I_l}\right),
		\qquad
		v_{S,y}^*|_{I_l}
		=
		\sum_{j=1}^{N_y}
		\omega_{y,j}
		\left(\lambda_{y,j}|_{I_l}\right),
		\\[6pt]
		p_D^*|_{I_l}
		=
		\sum_{k=1}^{N_{p_D}}
		\omega_{D,k}
		\left(\lambda_{D,k}|_{I_l}\right),
		\qquad l=1,\ldots,N_I.
	\end{gathered}
\end{equation}


Для более удобной сборки матриц перейдём к покомпонентной
записи граничных слагаемых.

Рассмотрим первое граничное слагаемое в уравнении Стокса
для компоненты скорости по $x$, подставляя разложение
давления Дарси, получаем
\begin{equation}
	\begin{gathered}
		\sum_{l=1}^{N_I}
		\int_{I_l}
		\left(
			\sum_{k=1}^{N_{p_D}}
			\omega_{D,k}\lambda_{D,k}|_{I_l}
		\right)
		n_{S,x}\Phi^{\mathcal I_\partial}_{i,l}
		\,d\Gamma
		=
		\sum_{k=1}^{N_{p_D}}
		\omega_{D,k}
		\left[
			\sum_{l=1}^{N_I}
			\int_{I_l}
			\left(\lambda_{D,k}|_{I_l}\right)
			n_{S,x}\Phi^{\mathcal I_\partial}_{i,l}
			\,d\Gamma
		\right].
	\end{gathered}
\end{equation}

Для второго граничного слагаемого уравнения Стокса
по $x$ получаем
\begin{equation}
	\begin{gathered}
		\sum_{l=1}^{N_I}
		\int_{I_l}
		\beta(\mathbf v_S^*\cdot\boldsymbol\tau)
		\tau_x\Phi^{\mathcal I_\partial}_{i,l}
		\,d\Gamma
		=
		\\[6pt]
		=
		\sum_{j=1}^{N_x}
		\omega_{x,j}
		\left[
			\sum_{l=1}^{N_I}
			\int_{I_l}
			\beta
			\left(\lambda_{x,j}|_{I_l}\right)
			\tau_x^2
			\Phi^{\mathcal I_\partial}_{i,l}
			\,d\Gamma
		\right]
		+
		\sum_{j=1}^{N_y}
		\omega_{y,j}
		\left[
			\sum_{l=1}^{N_I}
			\int_{I_l}
			\beta
			\left(\lambda_{y,j}|_{I_l}\right)
			\tau_y\tau_x
			\Phi^{\mathcal I_\partial}_{i,l}
			\,d\Gamma
		\right].
	\end{gathered}
\end{equation}

Аналогично преобразуются граничные слагаемые уравнения
Стокса для компоненты скорости по $y$.

Для граничного слагаемого уравнения Дарси получаем
\begin{equation}
	\begin{gathered}
		-
		\sum_{l=1}^{N_I}
		\int_{I_l}
		(\mathbf v_S^*\cdot\mathbf n_S)
		G^{\mathcal I_\partial}_{i,l}
		\,d\Gamma
		=
		\\[6pt]
		=
		-
		\sum_{j=1}^{N_x}
		\omega_{x,j}
		\left[
			\sum_{l=1}^{N_I}
			\int_{I_l}
			\left(\lambda_{x,j}|_{I_l}\right)
			n_{S,x}
			G^{\mathcal I_\partial}_{i,l}
			\,d\Gamma
		\right]
		-
		\sum_{j=1}^{N_y}
		\omega_{y,j}
		\left[
			\sum_{l=1}^{N_I}
			\int_{I_l}
			\left(\lambda_{y,j}|_{I_l}\right)
			n_{S,y}
			G^{\mathcal I_\partial}_{i,l}
			\,d\Gamma
		\right].
	\end{gathered}
\end{equation}


\newpage
Введём обозначения для коэффициентов граничных слагаемых
уравнения Стокса по $x$:
\begin{equation}
	\begin{gathered}
		(H_x)_{ik}
		=
		\sum_{l=1}^{N_I}
		\int_{I_l}
		\left(\lambda_{D,k}|_{I_l}\right)
		n_{S,x}\Phi^{\mathcal I_\partial}_{i,l}
		\,d\Gamma,
		\\[6pt]
		(R_{xx})_{ij}
		=
		\sum_{l=1}^{N_I}
		\int_{I_l}
		\beta
		\left(\lambda_{x,j}|_{I_l}\right)
		\tau_x^2\Phi^{\mathcal I_\partial}_{i,l}
		\,d\Gamma,
		\\[6pt]
		(R_{xy})_{ij}
		=
		\sum_{l=1}^{N_I}
		\int_{I_l}
		\beta
		\left(\lambda_{y,j}|_{I_l}\right)
		\tau_y\tau_x\Phi^{\mathcal I_\partial}_{i,l}
		\,d\Gamma.
	\end{gathered}
\end{equation}

Для коэффициентов граничного слагаемого уравнения Дарси
введём обозначения
\begin{equation}
	\begin{gathered}
		(C_x^D)_{ij}
		=
		\sum_{l=1}^{N_I}
		\int_{I_l}
		\left(\lambda_{x,j}|_{I_l}\right)
		n_{S,x}G^{\mathcal I_\partial}_{i,l}
		\,d\Gamma,
		\\[6pt]
		(C_y^D)_{ij}
		=
		\sum_{l=1}^{N_I}
		\int_{I_l}
		\left(\lambda_{y,j}|_{I_l}\right)
		n_{S,y}G^{\mathcal I_\partial}_{i,l}
		\,d\Gamma.
	\end{gathered}
\end{equation}


Подставляя разложение давления Дарси в объёмное слагаемое
и вынося коэффициенты за знаки интегралов, получаем
\begin{equation}
	\begin{gathered}
		\sum_{\alpha=1}^{N_D}
		\int_{\tau^D_\alpha}
		\frac{K}{\mu}
		\nabla\left(
			\sum_{k=1}^{N_{p_D}}
			\omega_{D,k}
			\lambda_{D,k}|_{\tau^D_\alpha}
		\right)
		\cdot\nabla G_{i,\alpha}
		\,d\tau^D_\alpha
		=
		\\[6pt]
		=
		\sum_{k=1}^{N_{p_D}}
		\omega_{D,k}
		\left[
			\sum_{\alpha=1}^{N_D}
			\int_{\tau^D_\alpha}
			\frac{K}{\mu}
			\nabla\left(
				\lambda_{D,k}|_{\tau^D_\alpha}
			\right)
			\cdot\nabla G_{i,\alpha}
			\,d\tau^D_\alpha
		\right].
	\end{gathered}
\end{equation}

Введём обозначение
\begin{equation}
	(A_D)_{ik}
	=
	\sum_{\alpha=1}^{N_D}
	\int_{\tau^D_\alpha}
	\frac{K}{\mu}
	\nabla\left(
		\lambda_{D,k}|_{\tau^D_\alpha}
	\right)
	\cdot\nabla G_{i,\alpha}
	\,d\tau^D_\alpha.
\end{equation}

Объединяя объёмные блоки Стокса с полученными блоками
границы сопряжения и области Дарси, запишем систему
в блочном виде:
\begin{equation}
	\begin{pmatrix}
		A_{xx}+R_{xx} & A_{xy}+R_{xy} & -B_x & H_x\\
		A_{yx}+R_{yx} & A_{yy}+R_{yy} & -B_y & H_y\\
		C_x & C_y & 0 & 0\\
		-C_x^D & -C_y^D & 0 & A_D
	\end{pmatrix}
	\begin{pmatrix}
		\mathcal W_x\\
		\mathcal W_y\\
		\mathcal W_S\\
		\mathcal W_D
	\end{pmatrix}
	=
	\begin{pmatrix}
		F_x\\
		F_y\\
		0\\
		0
	\end{pmatrix}.
	\label{eq:stokes_darcy_block_system}
\end{equation}



\subsection{Выбор тестовых и базисных функций}

Как и в разделе 4, воспользуемся методом Галёркина:
выберем тестовые функции совпадающими с базисными
функциями соответствующих искомых величин:
\begin{equation}
	\begin{gathered}
		\varphi_i=\lambda_{x,i},
		\qquad
		\psi_i=\lambda_{y,i},
		\\[6pt]
		q_i=\lambda_{S,i},
		\qquad
		g_i=\lambda_{D,i}.
	\end{gathered}
\end{equation}

Тогда число тестовых функций совпадает с числом базисных
функций соответствующего пространства:
\begin{equation}
	m_\varphi=N_x,\qquad
	m_\psi=N_y,\qquad
	m_q=N_{p_S},\qquad
	m_g=N_{p_D}.
\end{equation}

Для граничных блоков, связывающих скорость Стокса с
давлением Дарси, получаем
\begin{equation}
	C_x^D=H_x^T,
	\qquad
	C_y^D=H_y^T.
\end{equation}

Действительно, например, элементы $H_x$ и $C_x^D$
при выбранных тестовых функциях имеют вид
\begin{equation}
	\begin{gathered}
		(H_x)_{ik}
		=
		\sum_{l=1}^{N_I}
		\int_{I_l}
		\left(\lambda_{D,k}|_{I_l}\right)
		n_{S,x}
		\left(\lambda_{x,i}|_{I_l}\right)
		\,d\Gamma,
		\\[6pt]
		(C_x^D)_{ki}
		=
		\sum_{l=1}^{N_I}
		\int_{I_l}
		\left(\lambda_{x,i}|_{I_l}\right)
		n_{S,x}
		\left(\lambda_{D,k}|_{I_l}\right)
		\,d\Gamma
		=
		(H_x)_{ik}.
	\end{gathered}
\end{equation}
Для блоков $H_y$ и $C_y^D$ рассуждение аналогично.

Кроме того, матрица $A_D$ симметрична, а для блоков
касательного скольжения выполняются соотношения
\begin{equation}
	R_{xx}=R_{xx}^T,
	\qquad
	R_{yy}=R_{yy}^T,
	\qquad
	R_{xy}=R_{yx}^T.
\end{equation}


В системе~\eqref{eq:stokes_darcy_block_system} блоки,
связанные с давлением Стокса и давлением Дарси, стоят
с противоположными знаками в разных строках. Умножим
третье и четвёртое уравнения системы на $-1$. Тогда,
используя полученные соотношения, запишем
\begin{equation}
	\begin{pmatrix}
		A_{xx}+R_{xx} & A_{xy}+R_{xy} & -B_x & H_x\\
		A_{yx}+R_{yx} & A_{yy}+R_{yy} & -B_y & H_y\\
		-B_x^T & -B_y^T & 0 & 0\\
		H_x^T & H_y^T & 0 & -A_D
	\end{pmatrix}
	\begin{pmatrix}
		\mathcal W_x\\
		\mathcal W_y\\
		\mathcal W_S\\
		\mathcal W_D
	\end{pmatrix}
	=
	\begin{pmatrix}
		F_x\\
		F_y\\
		0\\
		0
	\end{pmatrix}.
\end{equation}

Объёмные блоки Стокса также удовлетворяют соотношениям
$A_{xx}=A_{xx}^T$, $A_{yy}=A_{yy}^T$ и
$A_{xy}=A_{yx}^T$. Следовательно, полученная матрица
является симметричной.

Для приближения давления в области Дарси выберем
непрерывные линейные конечные элементы $P_1$. На каждом
треугольнике области $\Omega_D$ определены три локальные
базисные функции, соответствующие его вершинам.
Эти функции строятся так же, как базисные функции давления
Стокса в разделе 4, и определяются по тем же формулам.
При их вычислении используются координаты вершин
соответствующего треугольника области Дарси.

\end{document}